\documentclass[11pt,xcolor={dvipsnames},hyperref={pdftex,pdfpagemode=UseNone,hidelinks,pdfdisplaydoctitle=true},usepdftitle=false]{beamer}
\usepackage{presentation}[aspectratio=169]
\usepackage{math}
\usepackage{mathtools}
\usepackage{mleftright}

% Enter title of presentation PDF:
\hypersetup{pdftitle={Aggregate Shocks, Taxes, and Debt Management II}}
% Enter link to PDF file with figures:


\begin{document}
% Enter presentation title:
\title{Aggregate Shocks, Taxes, and Debt Management II}
\subtitle{Fiscal and Monetary Policy 2023}
% Enter presentation information:

% Enter presentation authors:
\author{Piotr Żoch}%
% Enter presentation location and date (optional; comment line if not needed):
\frame{\titlepage}

% Fill out content of presentation:
\begin{frame}{Plan}
    \begin{itemize}
    \item We return to the questions posed last week and try to answer them in a different setting.
    \item We study the problem of a government that wants to maximize welfare of a representative household.
    \item The government can \al{fully commit} to its announced policies. This commitment is credible.
    \end{itemize}
\end{frame}


    \begin{frame}
        \heading{Lucas and Stokey (1983)}
    \end{frame}

    \begin{frame}{Setup}
    \begin{itemize}
    \item We study the following problem based on Lucas and Stokey (1983):
    \begin{itemize}
    \item the representative household has preferences over consumption and labor (leisure);
    \item production technology is linear in labor, there is no capital;
    \item the government uses \al{distortionary} tax rate $\tau_{t}$ on \al{labor income} to finance \emph{stochastic}
    purchases $g_{t}$ (we assume $g_{t}$ is a Markov process);
    \item the households and the government can trade a full set of \alg{Arrow securities} (state-contingent claims).
    \item the government wants to finance $g_{t}$ in a way that maximizes the welfare of the representative household.
    \end{itemize}
    \item General equilibrium.
    \end{itemize}
    \end{frame}


    %
    \begin{frame}{Setup}
    \begin{itemize}
    \item Price of securities no longer exogenous.
    \item The government's cares about welfare of the representative household.
    \item Think of the problem in two steps:
    \begin{enumerate} 
        \item for \al{given} government policies there exists some \alg{competitive equilibrium};
        \item the government picks policies that result in the "best" equilibrium.
    \end{enumerate}
    \item Key difference: policies not only have to satisfy the government's budget constraint, but also household's optimality conditions.
    \item \alb{Primal approach}: we will look directly for allocations that maximize welfare and then think of policies that implement that.
    \end{itemize}
    \end{frame}
    %
    \begin{frame}{Competitive equilibrium}
        \begin{beamerboxesrounded}[upper=block head,lower=block body,shadow=false]{\textbf{Competitive equilibrium}}
            A competitive equilibrium given government policies $\tau\bp{s^t}$ and $g\bp{s^t}$ is a set of allocations $\bp{c\bp{s^t},\ell\bp{s^t},a\bp{s^{t+1}}}$ and prices $\bp{q\bp{s^{t+1} \mid s^t },w\bp{s^t}}$ such that: 
            \begin{enumerate}
                \item Given prices, allocations solve the household problem;
                \item Given prices, allocations solve the firm problem;
                \item The government budget constraint is satisfied is each state;
                \item All markets clear in each state.
            \end{enumerate}
        \end{beamerboxesrounded}

        \begin{itemize}
            \item Note: this is not a proper definition of a competitive equilibrium, but it will do for our purposes.
        \end{itemize}
        \end{frame}
    
    \begin{frame}{Household problem}
    \begin{itemize}
    \item Household chooses consumption, labor supply and portfolio of Arrow securities to maximize expected utility:
    \item Household problem:
    \begin{align*}  
        &\max \E_0 \sum_{t=0}^\infty  \beta^t u\bp{c \bp{s^t},\ell\bp{s^t}}\\
        \text{s.t.} \quad \forall s^t \quad & c\bp{s^t}  + \sum_{s^{t+1} \geq  s^t} q\bp{s^{t+1}\mid s^t} a \bp{s^{t+1}} \\ &=\bp{1 - \tau\bp{s^t} }w \bp{s^t} \ell\bp{s^t} + a \bp{s^{t}}  \\
        & a \bp{s^0}  = a_0 
        \end{align*}
        where $u\bp{\cdot}$ is the utility function, $c\bp{\cdot}$ consumption and $\ell \bp{\cdot}$ labor.
          \end{itemize}
    \end{frame}


    \begin{frame}{Household problem}
    \begin{itemize}
    \item The sequence of household budget constraints can be written as a single constraint
    \begin{align*}  
    \sum_{t=0}^\infty \sum_{s^t}  q\bp{s^t} c\bp{s^t}  =
    \sum_{t=0}^\infty \sum_{s^t}  q\bp{s^t} \bp{1 - \tau\bp{s^t} }w \bp{s^t} \ell\bp{s^t} +q\bp{s^0} a_0 
    \end{align*}
    where $q\bp{s^t}$ is now the time-0 price of one unit of goods in state $s^t$.
    \item Why? \al{Complete markets} are \alg{equivalent} with Arrow securities and time-0 trading.
        \end{itemize}
    \end{frame}

    \begin{frame}{Household problem}
    \begin{itemize}
    \item Household problem:
    \begin{align*}  
        & \max \sum_{t=0}^\infty \sum_{s^t}  \beta^t \pi\bp{s^t} u\bp{c \bp{s^t},\ell\bp{s^t}}\\
        \text{s.t.} & \sum_{t=0}^\infty \sum_{s^t}  q\bp{s^t} c\bp{s^t}  =
        \sum_{t=0}^\infty \sum_{s^t}  q\bp{s^t} \bp{1 - \tau\bp{s^t} }w \bp{s^t} \ell\bp{s^t} \\ &  + q\bp{s^0}a_0 
        \end{align*}
        \item This is nice because there is only one constraint. 
        \item We could have used the same trick last week. 
            \end{itemize}
    \end{frame}


    %
    \begin{frame}{First order conditions}
    \begin{itemize}
    \item Write down the Lagrangian and differentiate with respect to $c\bp{s^t}$ and $\ell\bp{s^t}$ to get the first order conditions:
    \begin{align*}
    \bs{c\bp{s^t}:} \;\beta^t \pi\bp{s^t} u_c \bp{c \bp{s^t},\ell\bp{s^t}} & =\lambda q\bp{s^t}\\
    \bs{\ell\bp{s^t}:} \;\beta^t \pi\bp{s^t} u_{\ell} \bp{c \bp{s^t},\ell\bp{s^t}} & =-\lambda q\bp{s^t}\bp{1 - \tau\bp{s^t} }w \bp{s^t}
    \end{align*}
    \item Note: the Lagrange multiplier $\lambda$ is \alg{the same} for all states $s^t$.
    \item To simplify notation 
    \begin{align*}
        \bs{c\bp{s^t}:} \;\beta^t \pi\bp{s^t} u_c \bp{s^t} & =\lambda q\bp{s^t}\\
        \bs{\ell\bp{s^t}:} \;\beta^t \pi\bp{s^t} u_{\ell} \bp{s^t} & =-\lambda q\bp{s^t}\bp{1 - \tau\bp{s^t} }w \bp{s^t}
        \end{align*}
    \end{itemize}
    \end{frame}

\begin{frame}{Implementability constraint}
        \begin{itemize}
            \item The first order conditions:
        \begin{align*}
            \bs{c\bp{s^t}:} \;\beta^t \pi\bp{s^t} u_c \bp{s^t} & =\lambda q\bp{s^t}\\
            \bs{\ell\bp{s^t}:} \;\beta^t \pi\bp{s^t} u_{\ell} \bp{s^t} & =-\lambda q\bp{s^t}\bp{1 - \tau\bp{s^t} }w \bp{s^t}
        \end{align*}
            together with the budget constraint
        \begin{align*}
            \sum_{t=0}^\infty \sum_{s^t}  q\bp{s^t} c\bp{s^t}  =
            \sum_{t=0}^\infty \sum_{s^t}  q\bp{s^t} \bp{1 - \tau\bp{s^t} }w \bp{s^t} \ell\bp{s^t}  + q\bp{s^0}a_0   
        \end{align*}
        allow us to write the \alb{implementability constraint}:
         \begin{align*}  
            \sum_{t=0}^\infty \sum_{s^t} \beta^t \pi\bp{s^t} \bs{  u_c \bp{s^t} c\bp{s^t}  + u_{\ell} \bp{s^t} \ell\bp{s^t}}  = u_c \bp{s^0} a_0
        \end{align*}
        \end{itemize}
    \end{frame}
    \begin{frame}{Implementability constraint}
        \begin{itemize}
            \item The \alb{implementability constraint}
            \begin{align*}  
                \sum_{t=0}^\infty \sum_{s^t} \beta^t  \pi\bp{s^t} \bs{  u_c \bp{s^t} c\bp{s^t}  + u_{\ell} \bp{s^t} \ell\bp{s^t}}  = u_c \bp{s^0} a_0
            \end{align*}
        summarizes household optimality conditions and the budget constraint.
        \item The government can only choose allocations that satisfy this constraint.
        \item Captures the notion that the government's choices result in some \al{optimal} (given these choices) household behavior.
        \end{itemize}
    \end{frame}
    
    \begin{frame}{Firm problem}
        \begin{itemize}
            \item Output in this economy is produced by a representative price-taking firm.
            \item Linear production function that uses labor as an input.
            \item Firm problem is almost trivial -- it chooses labor to maximize profits:
            $$\max_{\ell\bp{s^t}} A\ell\bp{s^t} - w\bp{s^t} \ell \bp{s^t}.$$
            \item In a competitive equilibrium we must have
            $$ w\bp{s^t} = A.$$
        \end{itemize}
    \end{frame}
    \begin{frame}{Government budget constraint}
        \begin{itemize}
            \item The \alb{implementability constraint} is the household budget constraint (+ optimality conditions)
            \item The government should also be constrained by its budget constraint.
            \item We can ignore it and focus \al{directly} on the \al{resource constraint} of the economy: 
            $$ A \ell\bp{s^t} = g\bp{s^t} + c\bp{s^t}$$
            where $A \ell\bp{s^t}$ is the linear production technology.
            \item Why? \al{Walras' law} -- the government budget constraint is redundant.
        \end{itemize}
    \end{frame}
    


    \begin{frame}{Government problem}
    \begin{itemize}
    \item The government chooses \alg{allocations} $c\bp{s^t},\ell\bp{s^t}$ to maximize the household's utility 
    \begin{align*}
        \sum_{t=0}^\infty \sum_{s^t}  \beta^t \pi\bp{s^t} u\bp{c \bp{s^t},\ell\bp{s^t}}
    \end{align*}
    subject to the \alb{implementability constraint}:
    \begin{align*}  
        \sum_{t=0}^\infty \sum_{s^t}  \beta^t \pi\bp{s^t} \bs{  u_c \bp{s^t} c\bp{s^t}  + u_{\ell} \bp{s^t} \ell\bp{s^t}}  = u_c \bp{s^0} a_0
    \end{align*}
    and \al{resource constraints}: 
    \begin{align*}
    \forall s^t \quad   A \ell\bp{s^t} = g\bp{s^t} + c\bp{s^t}
    \end{align*}
    \end{itemize}
    \end{frame}
    %
\begin{frame}{Government problem}
    \begin{itemize}
    \item \alg{Full commitment}: the government announces its state-contingent plans at time 0 and cannot change them later.
    \item Everyone else trusts the government and knows it will stick to its plans.
    \end{itemize}   
    \end{frame}

    \begin{frame}{Optimality conditions}
    \begin{itemize}
    \item The \alg{first order conditions} for $s^t > s^0$ (we will return to $s^0$ later):
    \begin{align*}
      u_c\bp{s^t} + \eta \bs{ u_{cc} \bp{s^t} c\bp{s^t} + u_c \bp{s^t} + u_{c\ell} \bp{s^t} \ell\bp{s^t} }  & =  \mu \bp{s^t} \\
      u_\ell\bp{s^t} +\eta \bs{ u_{c\ell} \bp{s^t} c\bp{s^t} + u_\ell \bp{s^t} + u_{\ell\ell} \bp{s^t} \ell\bp{s^t} }  & =  - \mu \bp{s^t} A
    \end{align*}
    where $\eta$ is the Lagrange multiplier on the \alb{implementability constraint} and $\mu\bp{s^t}$ is the Lagrange multiplier on the \al{resource constraint}.
    \item There is also a resource constraint for each state: 
    \begin{align*}
        A \ell\bp{s^t} = g\bp{s^t} + c\bp{s^t}
    \end{align*}
    \item This gives us, for each state, three equations in three unknowns: $c\bp{s^t},\ell\bp{s^t},\mu\bp{s^t}$.
    \end{itemize}
    \end{frame}
    %
    \begin{frame}{Optimality conditions}
        \begin{itemize}
        \item We also have a fourth unknown: $\eta$ and one more equation: the \alb{implementability constraint}.
        \item Key: $\eta$ is the only link between different states.
        \item Once we know it, $c\bp{s^t},\ell\bp{s^t},\mu\bp{s^t}$ are determined for each $s^t$ from 2 first order conditions and the resource constraint.
        \item For two different states with the same $g\bp{s^t}$, the solution must be the same!
    \end{itemize}
        \end{frame}

    \begin{frame}{Implications}
        \begin{itemize}
        \item What are the implications for taxes?
        \item From household first order condition we have 
        \begin{align*}
            \bp{1-\tau\bp{s^t}} w\bp{s^t} = - \frac{u_\ell\bp{s^t}}{u_c\bp{s^t}}
        \end{align*}
        and linear production technology implies $w\bp{s^t} = A$ so
        \begin{align*}
            1-\tau\bp{s^t}  = - \frac{u_\ell\bp{c \bp{s^t},\ell\bp{s^t}}}{u_c\bp{c \bp{s^t},\ell\bp{s^t}}} \frac{1}{A}.
        \end{align*}
        \item We have solved for $c \bp{s^t},\ell\bp{s^t}$, we can get $\tau\bp{s^t}$!
        \end{itemize}
    \end{frame}
            
    \begin{frame}{Implications}
        \begin{itemize}
        \item Since for two different states with the same $g\bp{s^t}$ consumption and labor are the same -- taxes must be the same!
        \item \alg{Conclusion:} very different from Barro (1979) -- no history dependence.
        \item If government purchases are i.i.d., taxes are i.i.d. as well.
        \item Contrast with the \alb{ Barro (1979) case}: taxes were a random walk.
        \item Contrast with the \alr{complete markets case last week}: taxes were constant.
        \end{itemize}
    \end{frame}

    \begin{frame}{Implications}
        \begin{itemize}
        \item Why is the result different from the reduced form complete markets case?
        \item No full tax smoothing (unless with particular household preferences), because of a different objective.
        \item When $g_t$ is high, the government could want to \al{increase} labor supply to finance it without a large drop in consumption.
        \end{itemize}
    \end{frame}


    \begin{frame}{Example}
        \begin{itemize}
        \item Let $u\bp{\cdot} = \frac{c^{1-\gamma}}{1-\gamma} - \psi \ell $.
        \item Then $u_c = c^{-\gamma}$, $u_{cc} = -\gamma c^{-\gamma-1}$, $u_{\ell} = -\psi$, $u_{\ell\ell} = 0$, $u_{c\ell} = 0$.
        \item The first order conditions 
      \begin{align*}
            u_c\bp{s^t}  + \eta \bs{ u_{cc} \bp{s^t} c\bp{s^t} + u_c \bp{s^t} + u_{c\ell} \bp{s^t} \ell\bp{s^t} } & =  \mu \bp{s^t} \\
            u_\ell\bp{s^t}  + \eta \bs{ u_{c\ell} \bp{s^t} c\bp{s^t} + u_\ell \bp{s^t} + u_{\ell\ell} \bp{s^t} \ell\bp{s^t} } & =  - \mu \bp{s^t} A
          \end{align*}
          are
          \begin{align*}
            \bs{1+\eta\bp{1-\gamma}} c\bp{s^t}^{-\gamma} & =  \mu \bp{s^t} \\
            \bs{1+\eta}\psi & =   \mu \bp{s^t} A
          \end{align*}
            
        \end{itemize}
    \end{frame}

    \begin{frame}{Example}
        \begin{itemize}
        \item For each state $s^t>s^0$ we have:
        \begin{align*}
            \frac{1+\eta}{\bs{1+\eta\bp{1-\gamma}}A}\psi & =   c\bp{s^t}^{-\gamma}
        \end{align*}
        so consumption is constant.
        \item Taxes satisfy  
        \begin{align*}
        1-\tau\bp{s^t}  =  \frac{\psi}{c\bp{s^t}^{-\gamma}} \frac{1}{A}:
        \end{align*}
        since $c\bp{s^t}$ is constant, taxes are \al{constant} as well.
        \item With these preferences households supply labor fully elastically, fluctuations in $g_t$ are absorbed by changes in production.
        \end{itemize}
    \end{frame}
    
    \begin{frame}{Lessons}
        \begin{itemize}
        \item \al{Rembember:} the above result is \al{not} general.
        \item What is general is that here allocations and taxes are \al{history independent} -- they are fully determined by the \alg{current} state.
        \item What is the optimal path of debt? Similar logic as in our example last week: state-contingent debt that depends only on the next period state.
        \end{itemize}
    \end{frame}   

    \begin{frame}{Time-0}
        \begin{itemize}
        \item We ignored one important issue: \alg{first order conditions} are different for period 0. 
        \item This is because of the $u_c\bp{s^0} a_0$ term in the \alb{implementability constraint}.
        \item First order conditions for $s^0$ are:    
        \footnotesize
         \begin{align*}
            u_c\bp{s^0}+ \eta \bs{ u_{cc} \bp{s^0} c\bp{s^0} + u_c \bp{s^0} + u_{c\ell} \bp{s^0} \ell\bp{s^0} - u_{cc} \bp{s^0} a_0 }  & =  \mu \bp{s^0} \\
            u_\ell\bp{s^0}+\eta \bs{ u_{c\ell} \bp{s^0} c\bp{s^0} + u_\ell \bp{s^0} + u_{\ell\ell} \bp{s^0} \ell\bp{s^0} - u_{c\ell} \bp{s^0} a_0  }  & =  - \mu \bp{s^0} A
          \end{align*}
          \normalsize
        \item Time-0 allocations will be different from other periods, even with the same level of $g_t$.
    \end{itemize}
    \end{frame}   

    \begin{frame}{Time-0}
        \begin{itemize}
        \item The government has incentives to play with taxes in period 0 to change $q\bp{s^0}$ and thus affect the value of debt/assets.
        \begin{itemize}
            \item For example: it could try to make debt due at time 0 worthless.
            \item This is \alg{good}, because it allows the government to create \alg{lower} tax distortions in the future.
        \end{itemize}
       \item No similar effect in other periods, forward looking agents take it into account.
    \end{itemize}
    \end{frame}   


    \begin{frame}{Time inconsistency}
        \begin{itemize}
        \item \al{Time consistency problem}: if the government could re-optimize at future dates, period 1 would be like period 0. 
        \item Same logic applies to all other periods.
        \item But forward-looking agents would know the government is tempted to do it every period!
        \item This is actually the essence of Lucas and Stokey (1983).
            \begin{itemize}
                        \item Is there a way to design debt portfolio (maturities etc.) to ensure the government will \al{not be tempted} to re-optimize?
            \end{itemize}
    \end{itemize}
    \end{frame}   

\end{document}